\section{A typed quiver model}

The path-algebra language can be seen in a small model that also displays a genuine missing-packet obstruction.

Let the quiver have vertices $U,V$ and arrows
\[
R:U\to U,
\qquad
A:U\to V,
\qquad
B:V\to U,
\]
with the path-algebra relation
\[
BA=e_U,
\]
whose realization is the operator identity $BA=I_U$.
Realize
\[
U=\bbk^2,
\qquad
V=\bbk^3,
\]
by
\[
R=
\begin{pmatrix}0&-1\\1&0\end{pmatrix},
\qquad
A(u_1,u_2)=(u_1,u_2,0),
\qquad
B(v_1,v_2,v_3)=(v_1,v_2).
\]
Starting from the seed $e_1\in U$, the $R$-orbit spans $U$, and $A$ carries it onto the first two coordinates of $V$. Any path from $U$ to $V$ ends with $A$ and therefore lands in $\im A$. The third coordinate, which is $\ker B$, is unreachable from a seed supported only at $U$. Therefore the generated representation is
\[
U\oplus\Span\{(1,0,0),(0,1,0)\}
\subsetneq U\oplus V.
\]
Adding the seed $(0,0,1)\in V$ generates the full representation.

\begin{proposition}[Typed reachability certificate]
\label{prop:typed-reachability}
In the preceding quiver realization, no seed supported at $U$ can generate the full representation, whereas the two seeds
\[
e_1\in U,
\qquad
e_3\in V
\]
do generate it.
\end{proposition}

\begin{proof}
The first claim follows because every path that starts at $U$ and ends in $V$ has final arrow $A$, so its image lies in $\im A$. For the second, $e_1$ and $Re_1$ span $U$; their $A$-images span $\im A$; and $e_3$ spans $\ker B$, the missing complementary direction in $V$.
\end{proof}

This example shows why a typed grammar is more than a notational convenience. The obstruction lies in the subspace $\ker B\subset V$, which is unreachable from seeds supported at $U$, and is visible before any coordinate calculation.

\section{Worked algebraic model: an elliptic coordinate ring}

Let $\operatorname{char}\bbk\ne2$ and let
\[
A=\bbk[x,y]/(y^2-x^3+x).
\]
The affine equation $y^2=x^3-x$ has classical discriminant $64$, which is nonzero under this hypothesis; standard algebraic geometry then identifies its smooth projective completion as an elliptic curve \cite{silverman2009}. The Lean development checks the algebraic implication $2\ne0\Rightarrow64\ne0$, the quotient normal form, and the operator identities below. It does not formalize projective smoothness, genus, or the geometric identification of the discriminant. The packet computation itself uses only the displayed quotient-ring relation.
Every class has a unique normal form
\[
f(x,y)=a(x)+y b(x).
\]
Indeed, $y^2-(x^3-x)$ is monic of degree two in $y$, so division in $\bbk[x][y]$ gives a unique remainder of degree less than two in $y$. Concretely, every power $y^m$ with $m\ge2$ can be reduced using
\[
y^2=x^3-x.
\]

\subsection{One module seed and two address operators}

Take the seed
\[
s=1
\]
and let $M_x,M_y$ be multiplication by $x$ and $y$. Then
\[
M_x^iM_y^js=x^iy^j,
\]
so the algebra, regarded as a module over its multiplication algebra, is generated by one seed.

The two packet spaces are
\[
A_0=\bbk[x],
\qquad
A_1=y\bbk[x].
\]
Let
\[
P_0(a+yb)=a,
\qquad
P_1(a+yb)=yb.
\]
These are $\bbk$-linear pairwise annihilating idempotents and satisfy $P_0+P_1=I$. They are not $A$-linear module maps: multiplication by $y$ exchanges the two packets.

Identify the packet coefficients by
\[
A_0\oplus A_1\longrightarrow\bbk[x]\oplus\bbk[x],
\qquad
(a,yb)\longmapsto(a,b).
\]
In these ordered coefficient coordinates, with $M_x$ and $M_{x^3-x}$ denoting multiplication on $\bbk[x]$,
\[
{
M_x=
\begin{pmatrix}
M_x&0\\
0&M_x
\end{pmatrix},
\qquad
M_y=
\begin{pmatrix}
0&M_{x^3-x}\\
I&0
\end{pmatrix}.
}
\]
The relation is reconstructed as an operator identity:
\[
M_y^2=M_x^3-M_x.
\]

\begin{figure}[htbp]
\centering
\includegraphics[width=0.61\textwidth]{08_elliptic_packet_transfer.pdf}
\caption{Multiplication by $y$ exchanges the even and odd packets. The odd-to-even transfer carries the relation weight $x^3-x$.}
\label{fig:elliptic-transfer}
\end{figure}

\subsection{Exact product law}

Write
\[
f=a_0+ya_1,
\qquad
g=b_0+yb_1.
\]
Then
\[
{
fg
=
\bigl(a_0b_0+(x^3-x)a_1b_1\bigr)
+
y\bigl(a_0b_1+a_1b_0\bigr).
}
\]
This is a two-packet convolution law. It is already visible at the linear operator level because multiplication by a fixed element is a transfer operator.

\begin{example}[Worked reduction]
Take
\[
f=x^2+y,
\qquad
g=x-y.
\]
Then
\[
\begin{aligned}
fg
&=x^3-x^2y+xy-y^2\\
&=x^3-x^2y+xy-(x^3-x)\\
&=x+(x-x^2)y.
\end{aligned}
\]
Thus
\[
P_0(fg)=x,
\qquad
P_1(fg)=(x-x^2)y.
\]
The symbolic certificate in the companion code performs the same quotient reduction exactly.
\end{example}

\subsection{What this example establishes}

The coordinate ring illustrates four distinct parts of the calculus:
\begin{enumerate}[label=(\roman*)]
\item the unit is a module seed;
\item $M_x$ and $M_y$ are address operators;
\item the curve equation is an operator relation;
\item parity in $y$ gives exact packet analysis, transfer, and reconstruction.
\end{enumerate}
It does not imply that every algebra has a one-generator algebra presentation. It shows instead that unital algebras are cyclic as modules over their full multiplication algebra.

\section{Worked signal model: a bilateral shift}

Let $\mathbb{F}\in\{\R,\C\}$ and let
\[
X=\ell^2(\Z;\mathbb F)
\]
with the explicit integer-indexed family of bounded reindexing operators
\[
(S_mx)[n]=x[n-m]
\qquad(m,n\in\Z).
\]
Thus $S_0=I$, $S_{m+r}=S_mS_r$, and $S_1$ is the usual bilateral shift. We use the standard Hilbert-space and bounded-operator conventions of \cite{conway1990}.
The impulse $\delta_0$ satisfies
\[
S_m\delta_0=\delta_m
\]
for every $m\in\Z$. Therefore finite-support sequences lie in its algebraic orbit hull and are dense in $\ell^2(\Z)$.

\begin{proposition}[Shift seed]
\label{prop:shift-seed}
For $n\in\N$, let $G(n)$ mean that some $n$-element finite seed set has dense hull under the declared family $\{S_m:m\in\Z\}$. Then
\[
{
1=\min\{n\in\N:G(n)\}.
}
\]
Equivalently, one seed suffices and no empty seed family generates densely.
\end{proposition}

\begin{proof}
The orbit of $\delta_0$ contains every standard basis vector. Their finite linear span is dense, so $G(1)$ holds. The hull of the empty seed set is zero, while $\ell^2(\Z;\mathbb F)$ is nonzero, so $G(0)$ fails. Therefore $1$ is the least admissible finite cardinality.
\end{proof}

For $N\ge1$, define residue projectors
\[
(P_{N,j}x)[n]
=
\begin{cases}
x[n],&n\equiv j\pmod N,\\0,&\text{otherwise.}\end{cases}
\]
They form an orthogonal packet system and satisfy
\[
S_1P_{N,j}=P_{N,j+1}S_1,
\]
with indices modulo $N$. The shift is therefore a cyclic transfer between residue packets.

\begin{figure}[htbp]
\centering
\includegraphics[width=0.66\textwidth]{07_shift_packet_transfer.pdf}
\caption{Residue-packet support in the 24-coordinate periodic shift model with six packets. Each input residue transfers into the next output residue. The infinite bilateral-shift theorem is established separately in the text.}
\label{fig:shift-transfer}
\end{figure}

\begin{boundary}
This section uses residue packets only to demonstrate the foundational transfer algebra. Nested refinement and filter-bank structures require additional definitions and are outside Paper I.
\end{boundary}

\section{Worked matrix models}

\subsection{The full matrix algebra}

Let $X=\bbk^n$ and let $\cA=M_n(\bbk)=\End_{\bbk}(X)$, with the full matrix algebra declared before a seed or target is chosen. By \cref{cor:full-matrix}, every nonzero seed is cyclic. This is the maximal linear grammar on $X$, not a compression result; under a smaller declared operator algebra, the same vector need not be cyclic.

\subsection{Companion and repeated-block operators}

For one operator $T$, \cref{thm:single-operator} turns verified primary data
into an exact minimum-seed certificate.  The numerical companion compares:
\begin{enumerate}[label=(\alph*)]
\item a nilpotent companion operator with one Jordan chain, for which the Krylov rank is $1,2,\ldots,n$;
\item two equal Jordan blocks, for which one diagonal seed produces a chain only as long as one block and then stalls.
\end{enumerate}

The observed stall is not numerical failure.  The associated
two-dimensional scalar quotient is the lower-bound certificate in
\cref{thm:single-operator}; the identification of that quotient with the
geometric-multiplicity space is the classical Jordan input stated above.

\section{Worked graph model: adjacency algebras}

Let $G$ be a nonempty finite graph with adjacency matrix $A_G$, and regard
$\bbk^{V(G)}$ as a module over $\bbk[A_G]$.  The general result applies directly to a spectral certificate. For a symmetric real adjacency matrix, the classical spectral theorem supplies such data.

\begin{corollary}[Certified graph seed number]
\label{cor:graph-rank}
Let $\bbk$ be a field and let $\Lambda$ be a nonempty finite set.  Suppose
there are pairwise distinct scalars $a_{\lambda}\in\bbk$, multiplicities
$m_{\lambda}\in\N$, and an eigenbasis
\[
\{b_{\lambda,j}:\lambda\in\Lambda,\ 1\le j\le m_{\lambda}\}
\]
of $\bbk^{V(G)}$ satisfying
\[
A_Gb_{\lambda,j}=a_{\lambda}b_{\lambda,j}.
\]
If $m_{\lambda}\le R$ for all $\lambda$ and
$m_{\lambda_*}=R$ for some $\lambda_*$, then $R$ is the least natural number
of adjacency seeds:
\[
{
R=\min\left\{r\in\N:
\begin{array}{l}
\text{there exist }s_1,\ldots,s_r\text{ with}\\[2pt]
\bbk^{V(G)}=\Span_{\bbk}
  \{p(A_G)s_j:p\in\bbk[t],\ 1\le j\le r\}
\end{array}
\right\}.
}
\]
\end{corollary}

\begin{proof}
The distinct eigenvalue labels give Lagrange polynomials that project onto
the individual eigenspace packets.  Padding one basis vector from every
packet into each of $R$ seeds gives the upper bound.  Projection to the
$m_{\lambda}$ eigenspace coordinates is a surjective scalar quotient, so any
generating family has at least $m_{\lambda}$ members for every $\lambda$.
The label $\lambda_*$ gives the matching lower bound $R$.
\end{proof}

\begin{boundary}[Spectral-classification input]
For a finite undirected graph over $\C$, the real-symmetric spectral theorem
classically supplies a certificate of the displayed kind, and the numbers
$m_{\lambda}$ are the usual eigenvalue multiplicities; see
\cite{brouwerHaemers2012}.  The formal declaration verifies the
seed-number formula from the supplied eigenbasis and distinct eigenvalue
labels.  It does not derive diagonalizability from symmetry or implement an
eigensolver.
\end{boundary}

\subsection{Path graphs}

For every field $\bbk$ and every $n\ge1$, an endpoint of $P_n$ is a cyclic
vector for the adjacency operator.  Consequently one is the least finite
adjacency-seed number:
\[
{
1=\min\left\{r\in\N:
\exists s_1,\ldots,s_r,\quad
\bbk^n=\Span_{\bbk}\{p(A_{P_n})s_j:p\in\bbk[t],\ 1\le j\le r\}
\right\}.
}
\]
Over $\C$, the familiar spectral description has the distinct eigenvalues
\[
\lambda_k=2\cos\frac{k\pi}{n+1},
\qquad 1\le k\le n,
\]
and the normalized eigenvectors have endpoint coordinate proportional to
\[
\sin\frac{k\pi}{n+1}\ne0.
\]
There is also an all-field combinatorial proof.  Let
\[
K_n=[e_1,Ae_1,\ldots,A^{n-1}e_1].
\]
Its $(i,j)$ entry counts walks of length $j-1$ from the first vertex to vertex $i$. This entry is zero when $i>j$, while the diagonal entry is one because there is a unique shortest walk of length $i-1$ from vertex $1$ to vertex $i$. Hence $K_n$ is upper triangular with unit diagonal, so
\[
\det K_n=1
\qquad(n\ge1).
\]
Thus the endpoint is cyclic for every path size, not merely for the tested range.

\begin{figure}[htbp]
\centering
\includegraphics[width=0.75\textwidth]{05_path_graph_seed_rank.pdf}
\caption{The endpoint Krylov certificate gives every nonempty path one adjacency seed.}
\label{fig:path-rank}
\end{figure}

\subsection{Cycle graphs}

Two adjacent vertex vectors generate the adjacency module of every simple
cycle.  Over $\C$, a primitive-root Fourier quotient supplies a matching
two-dimensional scalar obstruction.  Hence, for every $n\ge3$, two is the
least finite adjacency-seed number:
\[
{
2=\min\left\{r\in\N:
\exists s_1,\ldots,s_r,\quad
\C^n=\Span_{\C}\{p(A_{C_n})s_j:p\in\C[t],\ 1\le j\le r\}
\right\}.
}
\]
To see the upper bound directly, index the vertices modulo $n$.
The identity $Ae_j=e_{j-1}+e_{j+1}$ recovers successive vertex vectors
from $e_0$ and $e_1$, over any field. For the lower bound over $\C$, let
$\zeta$ be a primitive $n$th root of unity and define
\[
q(x)=\left(\sum_{j=0}^{n-1}\zeta^j x_j,
             \sum_{j=0}^{n-1}\zeta^{-j}x_j\right).
\]
Then $q(Ax)=(\zeta+\zeta^{-1})q(x)$. The images of $e_0$ and $e_1$
are $(1,1)$ and $(\zeta,\zeta^{-1})$, which are independent because
$\zeta\ne\zeta^{-1}$ for $n\ge3$. Thus $q$ is a two-dimensional
scalar quotient, and \cref{prop:scalar-quotient} gives the
matching lower bound.

The usual spectral description writes the eigenvalues as
\[
\lambda_k=2\cos\frac{2\pi k}{n},
\qquad 0\le k<n.
\]
The values for $k$ and $n-k$ coincide. Apart from the simple characters at $k=0$ and, for even $n$, $k=n/2$, the eigenspaces occur in pairs. Therefore the maximum eigenspace multiplicity is two, although some eigenspaces have multiplicity one, and
the certified minimum above is two.

\begin{figure}[htbp]
\centering
\includegraphics[width=0.75\textwidth]{06_cycle_graph_seed_rank.pdf}
\caption{Cycle graphs have maximum adjacency-eigenspace multiplicity two and therefore require two adjacency seeds, although one or two exceptional eigenspaces are simple.}
\label{fig:cycle-rank}
\end{figure}

The path and cycle examples use nearly identical local operators.  Their
minimum finite seed numbers differ because the path has an
endpoint Krylov basis, whereas the cycle has a two-dimensional Fourier
quotient obstruction.

\section{The common operator language}

The preceding examples can be written without changing the foundational notation.

\begin{longtable}{L{0.12\textwidth}L{0.13\textwidth}L{0.19\textwidth}L{0.17\textwidth}L{0.16\textwidth}}
\caption{Foundational realizations of the common operator language.}\label{tab:dictionary}\\
\toprule
Realization & Seed & Operator grammar & Packets & Completion or quotient \\
\midrule
\endfirsthead
\toprule
Realization & Seed & Operator grammar & Packets & Completion or quotient \\
\midrule
\endhead
Coordinate ring & $1$ & $M_x,M_y$ & even/odd $y$ degree & quotient by $y^2-x^3+x$ \\
Shift system & $\delta_0$ & $(S_m)_{m\in\Z}$ & residue classes of $n$ & $\ell^2$ closure \\
Full matrix module & $s\ne0$ & $M_n(\bbk)$ & chosen idempotents & finite-dimensional \\
Single operator & $s_1,\ldots,s_r$ & $\bbk[T]$ & primary or spectral & finite-dimensional \\
Graph adjacency & vertex seeds & $\bbk[A_G]$ & eigenspaces & finite-dimensional \\
Typed grammar & vertex seeds & path algebra $\bbk Q/I$ & vertex idempotents & relation quotient \\
\bottomrule
\end{longtable}

The invariant content is:
\[
{
\begin{aligned}
&\text{generated hull: }\Hull_{\cA}(S),\\
&\text{seed complexity: }\srank_{\cA}(X),\\
&\text{packet coordinates: }x\leftrightarrow(P_{\alpha}x)_{\alpha},\\
&\text{operator coordinates: }T\leftrightarrow(P_{\beta}TP_{\alpha})_{\beta,\alpha},\\
&\text{relations: }F^m/N_S,\\
&\text{coordinate invariance: }P_{\alpha}\mapsto UP_{\alpha}U^{-1},\\
&\text{realization: }\Phi A=\varphi(A)\Phi.
\end{aligned}
}
\]

\begin{figure}[htbp]
\centering
\includegraphics[width=0.96\textwidth]{10_theorem_dependency.pdf}
\caption{Structural overview of the two proof families. The generator and packet arguments meet in the worked realizations. This diagram is not an exhaustive graph of Lean declaration dependencies.}
\label{fig:dependency}
\end{figure}

\section{Computational certificates}

The accompanying Python package is not used as a substitute for proof. It provides exact symbolic checks and finite numerical regression tests for formulas already proved. The computed ranges and random seeds are recorded explicitly. These checks do not extend a theorem beyond its hypotheses or certify an unrestricted numerical implementation. The journal companion also tests integer counts beyond floating-point precision and rejects the unsupported complex-symmetric spectral shortcut.

\subsection{Test suite}

The test suite checks:
\begin{enumerate}[label=(\roman*)]
\item projector idempotence, annihilation, and completeness;
\item exact vector analysis and synthesis;
\item exact operator reconstruction from all transfer blocks;
\item coordinate transport under a dense random similarity transform;
\item cyclicity of a companion operator;
\item the two-seed obstruction for repeated Jordan blocks;
\item the semisimple multiplicity formula in representative cases;
\item full orbit dimension for two representative nonzero seeds under the full matrix algebra;
\item exact reduction and packet multiplication in the elliptic coordinate ring;
\item seed ranks of path and cycle adjacency modules over the reported finite ranges;
\item the cyclic residue transfer pattern of a finite shift.
\end{enumerate}

\subsection{Recorded results}

\begin{table}[htbp]
\centering
\caption{Archived representative checks from the v1.3.1 journal-edition run. Fresh companion runs are reported separately in the release verification record.}
\label{tab:results}
\begin{tabular}{L{0.57\textwidth}r}
\toprule
Quantity & Result \\
\midrule
Projector idempotence residual & $0$ \\
Projector annihilation residual & $0$ \\
Projector completeness residual & $0$ \\
Operator reconstruction residual & $0$ \\
Transported-projector residual, maximum & $1.85\times10^{-15}$ \\
Companion seed rank & $1$ \\
Two equal Jordan blocks, seed rank & $2$ \\
$M_3(\C)$ acting on seven copies, seed rank & $3$ \\
Sampled path graph seed ranks, $2\le n\le12$ & $1$ \\
Sampled cycle graph seed ranks, $3\le n\le12$ & $2$ \\
Full $M_4(\R)$ orbit dimensions for two sampled nonzero seeds & $4,4$ \\
Automated tests & $33/33$ passed \\
\bottomrule
\end{tabular}
\end{table}

The finite shift pattern recovered by the experiment is
\[
\begin{pmatrix}
0&0&0&1\\
1&0&0&0\\
0&1&0&0\\
0&0&1&0
\end{pmatrix},
\]
which is the exact packet adjacency of a cyclic four-phase transfer.

\subsection{Algorithms}

The following finite constructions state mathematical specifications for the companion. Their Lean counterparts and exact scope are listed in the separate guide. They do not assert code extraction or a particular backend's performance.

\begin{construction}[Orbit-hull closure specification]
\label{con:orbit-hull-algorithm}
Let $V$ be finite-dimensional, let $(O_i)_{i\in I}$ be a finite operator
family, and let $S\subseteq V$.  Define
\[
W_0=\Span_{\bbk}S,
\qquad
W_{n+1}=W_n+\sum_{i\in I}O_i(W_n).
\]
Then there is an $n\in\N$ such that
\[
W_{n+1}=W_n
\qquad\text{and}\qquad
W_n=\Hull_{(O_i)}(S),
\]
where the right side is the least subspace containing $S$ and invariant under
every $O_i$.  Thus the ascending subspace iteration terminates and its stable
value is exactly the finite operator hull.
\end{construction}

\begin{construction}[Transfer-block table specification]
\label{con:transfer-block-algorithm}
Given finite projector systems $(P_i)_{i\in I}$ and $(Q_j)_{j\in J}$ and a
linear map $T$, form the table
\[
B_{ji}=Q_jTP_i.
\]
Its assembly reconstructs the input map exactly,
\[
\sum_{j\in J}\sum_{i\in I}B_{ji}=T,
\]
and its reported support relation is exactly
$\{(i,j):B_{ji}\ne0\}$.  The latter equality is the specification of the
transfer-graph output, not a claim that Lean extracts a decidable edge-list
routine for every coefficient field.
\end{construction}

\begin{construction}[Certified exact graph seed-rank specification]
\label{con:exact-graph-seed-rank}
For a finite square matrix $A$ over a field $K$, define the noncomputable exact
spectral quantity
\[
r_{\mathrm{spec}}(A)
=
\max_{\mu\in\Spec_K(A)}
\dim_K\ker(A-\mu I),
\]
with value zero for an empty spectrum.  Suppose an exact certificate supplies
a finite eigenbasis indexed by labels and multiplicities, the eigenvalue
equations, a multiplicity bound attained by one label, projector polynomials
that extract the padded eigenbasis generators, coverage of every spectral
value by a label, equality of the declared and actual eigenspace dimensions,
and membership of the maximizing eigenvalue in the spectrum.  Then
$r_{\mathrm{spec}}(A)$ is the least natural number of polynomial-orbit seeds
for $A$.

This is a verified certificate checker/specification.  It does not construct
the eigenbasis, prove that a symmetric matrix is diagonalizable, or implement
an exact eigensolver.  The companion Python code performs symbolic
computations in the reported finite graph range; the all-size path and cycle
results are proved separately above.
\end{construction}


\section{Reading guide: common confusions and diagnostic questions}

\subsection{Common confusions}
\begin{center}
\begin{longtable}{@{}>{\raggedright\arraybackslash}p{0.26\textwidth}>{\raggedright\arraybackslash}p{0.68\textwidth}@{}}
\toprule
\textbf{Do not conflate} & \textbf{Correct distinction}\\
\midrule
\endhead
Grammar with realized algebra & The grammar may contain relations or words that act as zero.  Seed rank depends only on the image action.\\
Seed with basis vector & A seed is a generator under an algebra action.  One seed may generate a high-dimensional module.\\
Cyclic vector with eigenvector & A cyclic vector must generate the whole module under operator words.  It is usually not an eigenvector.\\
Packet with orthogonal subspace & A packet is the range of an idempotent in a complete annihilating family.  Orthogonality is optional.\\
Exact reconstruction with good conditioning & Algebraic inverse identities may coexist with large analysis and synthesis norms.\\
Nonzero transfer edge with surjective transfer & A nonzero block can have small rank.  Graph reachability alone does not prove generation.\\
Packet-pure seed rank with unrestricted seed rank & Mixed-support seeds may meet several source components simultaneously.\\
Continuous surjection with quotient isomorphism & The target topology must agree with the quotient topology.\\
Repeated notation with semantic realization & A true realization must preserve a declared collection of operations and identities.\\
\bottomrule
\end{longtable}
\end{center}

\subsection{Diagnostic questions for a new application}
Before claiming that a new subject is a realization of Paper I, answer the following.
\begin{enumerate}
\item What is the fixed operator grammar or acting algebra?
\item What is the concrete realization space or module?
\item Which seeds are declared before the target is chosen?
\item Is generation algebraic, topological, probabilistic, or approximate?
\item What are the idempotent packets, and in what sense do they sum to the identity?
\item Are the packets finite, countable unconditional, or merely formal?
\item What are the transfer blocks, and which blocks are known to vanish?
\item Which morphisms preserve the packet system?
\item Which quotient, completion, tensor, or limit operations must be preserved?
\item What exact theorem from Paper I is being specialized, and what new domain theorem remains to be proved?
\end{enumerate}

\subsection{One-page theorem spine}
\ConceptChain{
\(
\begin{array}{c}
\text{free seed presentation}\\[0.25em]
\Downarrow\\[-0.1em]
\text{minimal generator number}\\[0.25em]
\Downarrow\\[-0.1em]
\text{Jacobson-radical reduction}\\[0.25em]
\Downarrow\\[-0.1em]
\text{semisimple or primary computation}
\end{array}
\qquad\qquad
\begin{array}{c}
\text{packet idempotents}\\[0.25em]
\Downarrow\\[-0.1em]
\text{exact analysis/synthesis}\\[0.25em]
\Downarrow\\[-0.1em]
\text{transfer-block algebra}\\[0.25em]
\Downarrow\\[-0.1em]
\text{graphs, quotients, and realization}
\end{array}
\)
}

